\chapter{Écrire des règles en trois fonctions}

Un moteur de règles complet tient en trois méthodes.

\begin{center}
\begin{tabular}{@{}p{3.6cm}p{9.8cm}@{}}
\toprule
\textbf{Méthode} & \textbf{La question à laquelle elle répond} \\
\midrule
\texttt{Construire} & à quoi la partie ressemble au départ \\
\texttt{CoupsPossibles} & ce qu'on a le droit de faire \\
\texttt{Appliquer} & ce qui se passe quand on le fait \\
\bottomrule
\end{tabular}
\end{center}

Tout le reste --- créer et détruire une instance, cloner un état, le sérialiser,
le hacher, dire si la partie est finie, dire qui a gagné, vérifier qu'un coup est
légal --- est \textbf{le même pour tout le monde}, et il est déjà écrit.

Le fichier de ce chapitre existe, il compile, et l'atelier le charge :
\nkfile{Applications/ConquerorLab/exemples/rules/RegleFacile.cpp}.

\section{Le compte, mesuré}

\texttt{NkcRulesVTable} compte \textbf{24 entrées}. Voici où elles vont.

\begin{nkcode}{ou vont les 24 entrees de la vtable}
24  entrees de la vtable
 3  contiennent VOTRE jeu    Construire, CoupsPossibles, Appliquer
18  sont ecrites une fois pour toutes dans ConquerorRegleFacile.h
 3  sont optionnelles (geometrie) et restent nulles
\end{nkcode}

Les recopier à chaque module, ce serait dix-huit occasions de se tromper sur du
code qui n'a aucun rapport avec le jeu qu'on essaie de concevoir.

\section{Le module entier}

Voici un moteur qui joue vraiment. Rien n'a été coupé.

\begin{nkcode}{exemples/rules/RegleFacile.cpp --- le module complet}
#include "Conqueror/ConquerorRegleFacile.h"

using namespace nkentseu;
using namespace nkentseu::conqueror;
using namespace nkentseu::conqueror::facile;

struct MesRegles {

    static constexpr int32 kCote = 5;

    // 1. A quoi la partie ressemble au depart.
    void Construire(Grille &g) {
        g.topologie = NkcTopology::Square4;
        g.nbJoueurs = 2;
        g.AjouterRectangle(kCote, kCote);
        g.PoserAuxCoins(2);
    }

    // 2. Ce qu'on a le droit de faire : dupliquer vers une case voisine VIDE.
    void CoupsPossibles(const Partie &p, ListeCoups &out) {
        NkcCoord voisins[8];
        for (int32 i = 0; i < p.nbCases; ++i) {
            if (p.cases[i].owner != static_cast<int8>(p.joueur)) continue;
            const NkcCoord de = p.ou[i];
            const int32    n  = p.Voisins(de, voisins, 8);
            for (int32 k = 0; k < n; ++k)
                if (p.Vide(voisins[k]))
                    out.Dupliquer(p.joueur, de, voisins[k]);
        }
        // PASSER est ajoute tout seul si cette liste reste vide.
    }

    // 3. Ce qui se passe quand on le fait.
    void Appliquer(Partie &p, const NkcMove &m, Evenements &ev) {
        p.Poser(m.to, m.player);
        ev.Duplique(m.player, m.from, m.to);
        p.conquete[m.player] += 10;   // en DIXIEMES : 10 == 1,0 point

        NkcCoord    voisins[8];
        const int32 n = p.Voisins(m.to, voisins, 8);
        for (int32 k = 0; k < n; ++k) {
            if (!p.Ennemie(voisins[k], m.player)) continue;
            const int8 ancien = p.Proprietaire(voisins[k]);
            p.Poser(voisins[k], m.player);
            ev.Retourne(m.player, voisins[k], ancien);
        }
        p.PasserLaMain();
    }
};

NKC_REGLES(MesRegles, "MesRegles", "1.0.0", "Moi")
\end{nkcode}

La dernière ligne fabrique les vingt et une entrées de la vtable et les deux
symboles exportés. \textbf{C'est tout le module.}

\section{L'essayer}

\begin{enumerate}
  \item copiez \texttt{exemples/rules/RegleFacile.cpp} vers
        \texttt{travail/rules/mes\_regles.cpp} ;
  \item changez le nom dans \texttt{NKC\_REGLES} --- sinon deux entrées
        identiques au menu ;
  \item sauvegardez, attendez une seconde ;
  \item panneau \textbf{Modules} $\rightarrow$ votre module apparaît
        $\rightarrow$ \textbf{Nouvelle partie}.
\end{enumerate}

Si la compilation échoue, la sortie complète du compilateur s'affiche dans le
panneau \textbf{Modules}. C'est votre seul retour : lisez-la.

\section{Ce que l'échafaudage vous coûte}

Il n'y a pas de magie, et la contrepartie se dit en une phrase : \textbf{\texttt{Partie}
est une structure fixe, plate, sans pointeur.}

C'est ce qui permet aux dix-huit fonctions d'être justes \textbf{par
construction} et non par relecture.

\begin{center}
\begin{tabular}{@{}p{5.2cm}p{8.2cm}@{}}
\toprule
\textbf{ce que le cadre doit faire} & \textbf{comment, parce que \texttt{Partie} est plate} \\
\midrule
cloner un état & une affectation \\
sérialiser & un \texttt{memcpy} \\
hacher & \textsc{fnv-1a} sur les octets \\
dire si un coup est légal & énumérer les coups possibles et comparer \\
\bottomrule
\end{tabular}
\end{center}

Le prix : vous ne pouvez pas ranger dans \texttt{Partie} un état de taille libre
--- une liste qui grandit, un cache, un arbre. Tant que votre jeu tient dans les
cases, les joueurs, l'énergie et la conquête, vous n'y penserez jamais.

\section{Le palier 1 : \textsc{fusionner} --- et comment ça se joue à l'écran}

C'est la question qui bloque le plus de monde, alors on y répond dans l'ordre :
\textbf{d'abord le geste}, ensuite le code.

\subsection*{Le geste, à la souris}

Exactement le même que \textsc{dupliquer} : \textbf{deux clics.}

\begin{enumerate}
  \item vous cliquez \textbf{une des cases à consommer} : elle prend le contour
        orange «~sélection~», et des \textbf{anneaux verts} apparaissent sur les
        destinations ;
  \item vous cliquez \textbf{l'anneau vert} : le coup part ;
  \item le totem résultat apparaît un niveau plus haut, le journal affiche
        \textbf{\texttt{FUSIONNER}}, et l'atelier joue l'animation de fusion.
\end{enumerate}

Vous \textbf{ne sélectionnez pas} les cases une par une. Votre module a déjà
énuméré les groupes possibles dans \texttt{CoupsPossibles} ; l'interface ne fait
que \textbf{retrouver} le coup que vous désignez par ces deux clics.

\begin{nkpiege}[Réparé le 2026-08-23 --- et il faut savoir pourquoi]
Un coup \textsc{dupliquer} porte sa source dans \texttt{from}. Un coup
\textsc{fusionner} laisse \texttt{from} \textbf{à zéro} et met les cases
consommées dans \texttt{fuseCells[0..fuseCount-1]}. Or l'atelier cherchait la
case cliquée dans \texttt{from}, et seulement là : \textbf{toutes} les fusions
lui semblaient partir de la case \texttt{(0,0)}. Conséquence mesurée : une IA
jouait vos fusions sans broncher, le journal les affichait, et \textbf{vous ne
pouviez pas les cliquer}. Si vous avez un kit antérieur au 23/08, c'est ce que
vous constatiez --- le défaut était dans l'atelier, pas dans vos règles.
\end{nkpiege}

\subsection*{Le code : une trentaine de lignes de plus}

\nkfile{Applications/ConquerorLab/exemples/rules/RegleFusion.cpp} est
\nkfile{RegleFacile.cpp} avec cette seule action en plus. Copiez-le dans
\texttt{travail/rules/} et jouez-le : c'est plus court à lire que ce chapitre.

\begin{nkcode}{exemples/rules/RegleFusion.cpp --- generer le coup}
const NkcCoord paire[2] = {a, voisin};
out.Fusionner(p.joueur, paire, 2, a, niveau + 1);
//            joueur    cases  nb  vers  niveau vise
\end{nkcode}

\texttt{Fusionner} existe pour \textbf{la même raison que \texttt{Dupliquer}} :
le contrat compare les coups \textbf{octet par octet}, un coup de fusion a deux
champs de plus (\texttt{fuseCount}, \texttt{fuseCells}), et les cases
inutilisées du tableau doivent être à zéro. Sinon le coup que votre IA propose
n'est \textbf{pas égal} au coup que \texttt{CoupsPossibles} avait généré,
l'atelier le refuse, et rien dans votre logique n'est faux. Le \texttt{memset}
est dans le raccourci.

\begin{nkcode}{exemples/rules/RegleFusion.cpp --- l'appliquer}
if (m.kind == NkcMoveKind::Fuse) {
    for (int32 k = 0; k < m.fuseCount; ++k) p.Vider(m.fuseCells[k]);  // VIDER d'abord
    const int8 niveau = (m.targetLevel >= 0) ? m.targetLevel : 1;
    p.Poser(m.to, m.player, niveau);                                  // POSER ensuite
    ev.Fusionne(m.player, m.to, niveau);
    p.conquete[m.player] += 10 * niveau;
    p.PasserLaMain();
    return;
}
\end{nkcode}

\textbf{Trois pièges, dans l'ordre où ils vous tomberont dessus.}

\begin{enumerate}
  \item \textbf{Vider avant de poser.} \texttt{m.to} fait partie du groupe
        consommé. Poser avant de vider efface le totem qu'on vient de créer ---
        et ça ne se voit qu'à la dixième partie, quand un totem disparaît sans
        raison.
  \item \textbf{Ne proposer chaque paire qu'une fois.} Sans le filtre
        \texttt{Index(voisin) > i}, la paire \{A,B\} sort deux fois --- vue de
        A, puis vue de B. Rien ne plante : la distribution des coups est
        seulement fausse, et votre campagne mesure autre chose que ce que vous
        croyez.
  \item \textbf{Payer la fusion en conquête.} Un totem de niveau $n$ vaut $n+1$
        totems. Si la conquête ne le reflète pas, fusionner fait \textbf{perdre}
        des points, aucune IA ne le fera jamais, et vous chercherez le bug dans
        votre générateur de coups.
\end{enumerate}

\begin{nkretenir}[Ce qu'on en mesure]
Deux campagnes, \textbf{même graine}, fusion activée puis désactivée. Si la
durée des parties ne bouge pas, votre fusion ne sert à rien : le problème est
dans le \textbf{coût} que vous lui donnez, pas dans le code.

\begin{nkretenir}[Comment se joue une fusion à la souris ?]
Ce n'est pas dans ce chapitre : il est écrit du point de vue de celui qui
\textbf{écrit} le moteur. Le geste, les marques à l'écran et le menu
«~Quel coup ?~» sont au chapitre \textbf{«~Jouer un coup à la souris~»}.
\end{nkretenir}

Mesure du 2026-08-23 sur \nkfile{RegleFusion.cpp} (carré 5$\times$5, 2 joueurs,
en prenant systématiquement la fusion quand elle existe) : \textbf{196 coups
joués dont 87 fusions, niveau 3 atteint, tous les coups légaux, rejeu
déterministe.}
\end{nkretenir}

\section{Artefacts et pouvoirs : ce qui existe, et ce qui n'existe pas}

Vous voulez anticiper --- bien. Voici l'état \textbf{exact}, pour ne pas
construire sur du vide.

\begin{center}
\begin{tabular}{@{}p{2.4cm}p{3.6cm}p{3.6cm}p{3.4cm}@{}}
\textbf{} & \textbf{Dans le contrat} & \textbf{Dans l'atelier} & \textbf{Dans un module livré} \\
\textsc{fusionner} & oui --- \texttt{NkcMoveKind::Fuse}, \texttt{fuseCells}, \texttt{fuseCount}, \texttt{targetLevel} & oui --- jouable à la souris depuis le 23/08, journal \texttt{FUSIONNER} & \textbf{oui} --- \texttt{exemples/rules/RegleFusion.cpp} \\
\textsc{pouvoir} & oui --- \texttt{NkcMoveKind::Power}, \texttt{powerId}, \texttt{NkcCellView::power} & \textbf{oui depuis le 23/08} --- clic sur le lanceur, clic sur la cible ; étiquette \texttt{P0}, \texttt{P1}, menu «~Quel coup ?~» si plusieurs pouvoirs visent la même case & non \\
\textsc{artefact} & oui --- \texttt{NkcCellView::artefact}, \texttt{ArtefactPlaced} / \texttt{ArtefactExpired} & \textbf{non} --- rien ne les dessine & non \\
\end{tabular}
\end{center}

Autrement dit : \textbf{les trois sont prévus dans le contrat, deux sont jouables
à la souris, aucun module livré ne pose encore d'artefact.}

\begin{itemize}
  \item Un \textbf{pouvoir} est un coup \texttt{Power} avec \texttt{from} = le
        lanceur, \texttt{to} = la cible, \texttt{powerId} = lequel. Le geste est
        donc le même que pour dupliquer : \textbf{clic sur le lanceur, clic sur
        la cible}. Le troisième champ, \texttt{powerId}, est celui que deux clics
        ne peuvent pas dire --- c'est pourquoi l'atelier ouvre le menu
        \textbf{«~Quel coup ?~»} dès que deux coups partagent la même
        destination. Avant le 23/08 il jouait « le premier coup qui correspond » :
        le second pouvoir d'un totem sur une même cible était
        \textbf{injouable à la souris}, sans le moindre message. Le raccourci
        \texttt{Pouvoir(joueur, de, vers, idPouvoir)} construit le coup avec la
        mise à zéro qu'exige le contrat. Mesurez tout de même vos pouvoirs
        \textbf{en campagne IA contre IA} : c'est là que se tranche
        l'équilibrage.
  \item Un \textbf{artefact} est un état \textbf{de case}, pas un coup :
        \texttt{artefact} porte son id, et vous émettez
        \texttt{ArtefactPlaced} / \texttt{ArtefactExpired} quand il apparaît et
        disparaît. Le champ est \textbf{sérialisé} --- un rejeu reste donc
        valide --- mais rien ne l'affiche encore.
        \nkfile{ConquerorRulesV2.cpp} le remet à \texttt{-1} et s'arrête là :
        c'est du palier 2.
  \item \texttt{NkcCellView::powerUsed} existe pour la \textbf{réserve} : un
        pouvoir attaché à un totem, consommable une fois. Le contrat le
        transporte ; personne ne s'en sert encore.
\end{itemize}

\begin{nkretenir}[La règle, si vous vous lancez]
Ne codez pas contre un champ que rien ne lit. Écrivez le pouvoir, faites-le jouer
par l'IA, mesurez-le en campagne, et \textbf{dites dans votre fiche que vous
n'avez pas pu le vérifier à la main} --- c'est une réponse recevable, l'invention
ne l'est pas.
\end{nkretenir}

\section{Quand partir --- et ce n'est pas un échec}

\begin{quote}
\textbf{C'est un échafaudage, pas une cage.}
\end{quote}

Le jour où votre état ne rentre plus dans \texttt{Partie}, vous écrivez le
contrat nu. Les vingt-quatre entrées sont documentées au chapitre
\textbf{«~Quand l'échafaudage ne suffit plus~»}, à la fin de ce cours --- et il
se lira beaucoup mieux à ce moment-là, parce que vous aurez déjà vu les dix-huit
fonctions à l'œuvre.

Vous n'avez rien à défaire pour passer de l'un à l'autre : les deux produisent le
même module, avec les mêmes deux symboles. \nkfile{RegleFacile.cpp} et
\nkfile{RegleContratNu.cpp} jouent \textbf{exactement le même jeu} ---
comparez-les, c'est l'exercice le plus instructif du cours.

\section*{Exercices}

\begin{enumerate}
  \item \textbf{Portée 2.} Autorisez la duplication vers une case à deux pas.
        Une seule ligne change dans \texttt{CoupsPossibles} --- laquelle ?
  \item \textbf{Le déplacement.} Ajoutez l'action \textsc{deplacer} (la case de
        départ se vide). Que devient le décompte des totems ?
  \item \textbf{Hexagone.} Passez \texttt{g.topologie} à
        \texttt{NkcTopology::HexPointy}. Combien de voisins avez-vous
        maintenant, et quelle ligne de votre code l'a supposé ?
  \item \textbf{La cascade.} Émettez \texttt{ev.Cascade(...)} quand un coup
        retourne au moins deux ennemis, et regardez ce que l'atelier en fait à
        l'écran.
\end{enumerate}
